% !TeX program = xelatex
\documentclass[11pt,lang=zh]{impeart}

% 使用標準文件類別的 lang=zh 選項；中文用 Shanggu，拉丁文字用 Libertinus。
\UseFont{libertinus}[global]
\setmonofont[
  BoldFont=lmmonolt10-bold.otf,
  ItalicFont=lmmono10-italic.otf,
  BoldItalicFont=lmmonolt10-boldoblique.otf
]{lmmono10-regular.otf}
\UseFeatures{hyperlinks,tables}

\usepackage{catchfile}
\usepackage{metalogo}
\usepackage{pdfpages}

\newcommand{\IMPEVersionFile}{../../VERSION}
\newcommand{\IMPEShowcaseFile}{../showcase/impe-showcase.pdf}

\IfFileExists{\IMPEVersionFile}
  {\CatchFileEdef{\IMPEVersion}{\IMPEVersionFile}{\endlinechar=-1\relax}}
  {\PackageError{impe-manual}
     {缺少必要的程式碼倉庫資源 ../../VERSION}
     {請從 manual/zh-tw 編譯，或在程式碼倉庫根目錄執行 scripts/build_manual.ps1。}}

% SOURCE_DATE_EPOCH 固定時，保持生成的手冊可重現。
\ifdefined\XeTeXversion
  \special{pdf:trailerid [
    <f2f1c967f3905e268855d68445d8351b>
    <031bcc35a4934255e3bb21323efc4efc>
  ]}
\else\ifdefined\pdftrailerid
  \expandafter\pdftrailerid\expandafter{IMPE-\IMPEVersion-manual-zh-tw}
\fi
\fi

\title{IMPE \LaTeX{} System}
\subtitle{使用者手冊}
\author{YANG Sikai}
\date{版本 \IMPEVersion}

\newcommand{\IMPERepo}{https://github.com/KelvinYangBK67/IMPE-LaTeX-System}
\newcommand{\IMPEIssues}{https://github.com/KelvinYangBK67/IMPE-LaTeX-System/issues}
\newcommand{\IMPEShowcaseURL}{https://github.com/KelvinYangBK67/IMPE-LaTeX-System/blob/main/manual/showcase/impe-showcase.pdf}
\newcommand{\IMPETaggedShowcaseURL}{https://github.com/KelvinYangBK67/IMPE-LaTeX-System/blob/v\IMPEVersion/manual/showcase/impe-showcase.pdf}

\begin{document}

\maketitle

\begin{center}
\small 繁體中文翻譯版；技術內容以英文參考版為準
\end{center}

\tableofcontents

\section{導論}

IMPE \LaTeX{} System 是一套模組化 \XeLaTeX{} 框架，提供可重用的版面、字體路由、
多語排版與選用文件功能。

IMPE 由三個可以獨立設定的部分組成：
\emph{版面}、\emph{字體}與\emph{功能}。版面描述文件主要的視覺組織方式；字體族提供
全域或局部的字體行為，必要時也包括依書寫系統進行的路由；功能模組則提供數學、
超連結、引用、表格或圖片等選用功能。這些部分既可分別選擇，也可透過統一的範本介面
組合使用。

IMPE 的一般用法刻意保持簡短。包裝類別可以先建立合理的文件版面與預設全域字體設定，
之後使用者再載入文件實際需要的額外字體族或功能。若希望更明確地控制設定，也可在
單一的 \verb|\UseTemplateSet| 呼叫中宣告相同的各個組件。

\XeLaTeX{} 是 IMPE 的主要文件引擎。這項選擇使系統能依賴 Unicode 文字、
OpenType 字體、\texttt{fontspec} 以及 \XeTeX{} 的字元間機制，同時保留熟悉的
\LaTeX{} 文件模型。某些特殊的內部路徑在特定文字或渲染方式需要時可能會呼叫其他
引擎，但一般 IMPE 文件預期以 \XeLaTeX{} 編譯。

本手冊面向一般使用者。更詳細的子系統文件另行維護在程式碼
倉庫的 \texttt{docs/} 目錄下。尤其是字體、版面、功能與系統文件，記錄了註冊欄位與
實作細節；若在此重複，反而會使手冊變得冗長。英文手冊是參考版；各語言翻譯版應遵循
其結構與技術內容。

本系統的視覺概覽可參見專案的
\href{\IMPEShowcaseURL}{目前 IMPE 展示冊 PDF}。與本手冊一同凍結的版本可由
\href{\IMPETaggedShowcaseURL}{對應版本標籤}取得。這份標準展示冊也完整收入
附錄~\ref{app:showcase}，用於展示促成本系統設計的各類文字、書寫方向、版面與功能
組合。

\section{安裝}

\subsection{發行形式}

IMPE 1.1.0 準備了 \texttt{Full}、\texttt{Core} 與
\texttt{CTAN} 三種發行形式。三者使用相同的核心
\XeLaTeX{} 執行環境與標準文件類別，差別在於隨附字體及安裝程序。

\begin{description}
  \item[\texttt{Full}（完整包）]
    包含核心、相容入口、XSR，以及經授權檢查且實際收入發行包的
    可再分發字體檔。
  \item[\texttt{Core}（核心包）]
    包含核心、相容入口、XSR 與安裝工具。所需字體可由
    \TeX{} Live、作業系統或使用者提供的字體檔解析。
  \item[\texttt{CTAN}]
    根目錄為 \texttt{impe-framework/} 的
    \texttt{impe-framework.zip}，包含標準文件類別、
    套件程式碼、XSR 原始碼及文件，但不附帶字體檔。
    透過 CTAN 安裝不會自動執行 Python 套件安裝。
\end{description}

完整包實際包含哪些字體，依發行時經授權審核的清單決定。
字體登錄表中的所有項目不一定都隨包附送。

\subsection{安裝完整包或核心包}

直接安裝會將套件放入使用者 TEXMF 樹的
\texttt{TEXMFHOME/tex/latex/impe/}。標準公開入口為
\texttt{impe.sty}、\texttt{impeart.cls}、
\texttt{impereport.cls}、\texttt{impebook.cls}
和 \texttt{impebeamer.cls}。
舊式 \texttt{\_zh}、\texttt{next*} 名稱屬相容入口。

\begin{verbatim}
# Windows：已解壓的發行包
.\install.bat
# Windows：原始碼倉庫
.\scripts\install.ps1

# Unix：已解壓的發行包
sh ./install.sh
# Unix：原始碼倉庫
sh ./scripts/install.sh
\end{verbatim}

Windows 安裝程式可使用 \texttt{-TexmfRoot} 和
\texttt{-SkipXsrPython}；Unix 安裝程式可接受自訂 TEXMF 路徑。
若安裝時未能更新 \TeX{} 檔名資料庫，可另外執行
\texttt{mktexlsr}。透過 \TeX{} 發行版的套件管理器安裝時，
則遵循該發行版的安裝流程。

\subsection{XSR 與 Python 執行環境}

IMPE 隨附 XSR \textbf{1.3}。完整包與核心包的安裝程式
通常會為 XSR 建立獨立的 Python 虛擬環境，不會直接修改
使用者的全域 Python 套件。動態埃及聖書字排版、契丹小字堆疊、
圖片字形處理可能需要 Python \textbf{3.11 以上}、
\texttt{venv} 與 \texttt{pip}；首次安裝相依套件時也可能需要網路。

如果只要安裝 \TeX{} 部分，可以略過 Python 環境設定：

\begin{verbatim}
# Windows
.\scripts\install.ps1 -SkipXsrPython
# Unix
IMPE_SKIP_XSR_PYTHON=1 sh ./scripts/install.sh
\end{verbatim}

使用動態 XSR 功能時，需要允許 XeLaTeX 執行
Shell Escape，或先行預處理 XSR 內容。渲染程式還必須能讀取
實際 TTF／OTF 字體檔；只有系統字體名稱不足以讓 Python
取得字形輪廓。

\subsection{字體資源與本機設定}

IMPE 會先嘗試已登錄的本機字體檔，再向 \TeX{} 字體系統
或作業系統尋找字體。已能直接解析的系統字體無須另外複製
到 IMPE 目錄。新增字體檔通常放在
\texttt{assets/fonts/} 下，亦可在 \texttt{impe.local.tex}
內設定其他存放位置：

\begin{verbatim}
\SetCatalogFontRoot{D:/Fonts/IMPE}
\end{verbatim}

此命令也可在單一文件的導言區使用。
一般字體解析失敗時，IMPE 可以警告並採用代用字體；
XSR 則須取得所用字體的實體檔案。

\subsection{快速安裝檢查}

若系統提供 Libertinus，可用下面的最小文件確認類別、
字體登錄與局部命令能正常工作：

\begin{verbatim}
\documentclass{impeart}
\UseFont{libertinus}
\begin{document}
IMPE 已可使用。\LIB{Libertinus 局部文字。}
\end{document}
\end{verbatim}

這項檢查針對一般 \XeLaTeX{} 編譯；XSR 和發行包另有測試。

\section{快速開始}

最簡單的 IMPE 工作流程，是選擇一個包裝類別，載入所需的字體族或功能，然後照常撰寫
文件。IMPE 在一般 \LaTeX{} 文件結構之外提供可重用的設定層。

因此，一份簡短文章可以寫成：

\begin{verbatim}
\documentclass[11pt]{impeart}

\UseFont{libertinus}
\UseFeatures{math,hyperlinks}

\title{A Sample Document}
\subtitle{Typeset with IMPE}
\author{Author Name}

\begin{document}

\maketitle

\section{Introduction}
\LIB{Hello, world.}

\end{document}
\end{verbatim}

這個類別已經提供一般的英文文章版面與基礎字體。省略模式的
\verb|\UseFont{libertinus}| 會遵循已登錄的自動行為：將 Libertinus 登錄為明確的全域拉丁字體候選，
同時提供局部命令 \verb|\LIB|；\verb|\UseFeatures| 則啟用兩個選用功能模組。
其餘部分都是一般的 \LaTeX{}。

若某份文件更適合把設定集中成單一宣告，同樣的做法也可以寫成：

\begin{verbatim}
\documentclass{impeart}

\UseTemplateSet{
  fonts = {libertinus},
  features = {math,hyperlinks}
}
\end{verbatim}

範本組也可以選擇版面與其他字體族。實際使用時，包裝類別通常已提供合適的預設版面，
所以大多數文件可寫出與預設值不同的設定。

IMPE 也可以作為套件載入到標準 \LaTeX{} 類別之下：

\begin{verbatim}
\documentclass{article}
\usepackage{impe}

\UseTemplateSet{
  layout = en_doc,
  fonts = {libertinus},
  features = {hyperlinks}
}
\end{verbatim}

當底層類別必須明確保留，或希望手動選擇每一個 IMPE 組件時，這種形式特別合適。
包裝類別與直接載入套件使用的是同一套底層系統。

\section{可直接使用的範本}

程式碼倉庫在 \texttt{templates/} 下提供一小組起始專案。這些範本面向實際寫作，
回歸測試與內部展示則放在 \texttt{tests/}。

目前提供的起始範本包括：

\begin{description}
  \item[\texttt{article\_en/}] 使用 \texttt{impeart} 的英文文章起始範本。
  \item[\texttt{article\_zh/}] 使用 \texttt{impeart} 與 \texttt{lang=zh} 的中文文章範本。
  \item[\texttt{book\_en/}] 使用 \texttt{impebook} 的英文書籍起始範本。
  \item[\texttt{book\_zh/}] 使用 \texttt{impebook} 與 \texttt{lang=zh} 的中文書籍範本。
  \item[\texttt{beamer\_en/}] 使用 \texttt{impebeamer} 的英文簡報起始範本。
  \item[\texttt{beamer\_zh/}] 使用 \texttt{impebeamer} 與 \texttt{lang=zh} 的中文簡報範本。
\end{description}

每個範本都遵循已安裝套件的正常寫法，並包含實際可用的中繼資料、章節與正文結構；
在適合的情況下，也附有具代表性的表格、圖片或投影片內容。因此，新專案可以直接複製
最接近的範本目錄、修改正文，並沿用隨附的導言區。

這些範本刻意使用標準的 \texttt{impe*} 公開入口，代表一般使用者應採用的常規
寫法。僅供開發使用的 fixture 位於 \texttt{tests/}。

\section{文件類別與公開入口}

\subsection{標準文件類別與語言選擇}

新文件統一使用四個標準類別：\texttt{impeart}、
\texttt{impereport}、\texttt{impebook} 和
\texttt{impebeamer}。同一類別可以透過
\texttt{lang=en}（預設）或 \texttt{lang=zh}
切換語言與對應排版，無須另換類別名稱：

\begin{verbatim}
% 英文文章
\documentclass[12pt,twoside]{impeart}

% 中文文章
\documentclass[lang=zh,12pt]{impeart}

% 中文書籍
\documentclass[lang=zh]{impebook}

% 中文簡報
\documentclass[lang=zh]{impebeamer}
\end{verbatim}

這四個類別分別建立在 \texttt{article}、
\texttt{report}、\texttt{book} 與
\texttt{beamer} 之上。英文預設版面依次為
\texttt{en\_doc}、\texttt{en\_report}、
\texttt{en\_book}、\texttt{beamer}；
中文則為 \texttt{zh\_doc}、
\texttt{zh\_report}、\texttt{zh\_book}、
\texttt{beamer\_zh}。英文預設字體包含 CMU；
中文另加 Shanggu。使用者明確指定的字體可覆寫類別預設。

一般類別選項，如 \texttt{12pt}、\texttt{twoside}
及 Beamer 的 \texttt{aspectratio=169}，
會傳給底層標準類別。舊式 \texttt{impe*\_zh}
與 \texttt{next*} 類別仍可使用，但屬相容入口；
詳見第~\ref{sec:compatibility} 節。

\subsection{直接使用 \texttt{impe.sty}}

各包裝類別共用同一底層實作，並提供便利預設。使用者也可以保留標準類別，然後載入

\begin{verbatim}
\usepackage{impe}
\end{verbatim}

再明確指定版面、字體與功能。當文件類別由期刊、出版社或其他套件提供，且必須保留
該類別時，建議採用此方式。

\subsection{標題與副標題}

除了標準的 \verb|\title| 與 \verb|\author| 之外，IMPE 還支援
\verb|\subtitle{...}|。包裝類別會套用自己的標題區塊預設，但一般 \LaTeX{}
標題機制仍保持熟悉的形式。典型的導言區可以簡單寫成

\begin{verbatim}
\title{Main Title}
\subtitle{A More Specific Description}
\author{Author Name}
\end{verbatim}

並在正文中呼叫 \verb|\maketitle|。

\section{範本模型}

\subsection{四個設定鍵}

主要的統一設定命令是 \verb|\UseTemplateSet|。它公開的設定模型刻意保持簡單：

\begin{verbatim}
\UseTemplateSet{
  layout = <preset>,
  fonts = {<families>},
  globalfonts = {<families>},
  features = {<features>}
}
\end{verbatim}

四個鍵各有職責。

\begin{description}
  \item[\texttt{layout}] 選擇與目前文件類別相容的公開版面預設。
  \item[\texttt{fonts}] 依字體族已登錄的自動行為載入字體。這是一般情況下的正常
    選擇：一個字體族可以建立局部命令、套用文件全域字體，或安裝指定範圍的路由。
  \item[\texttt{globalfonts}] 明確強制列出的字體族使用全域模式或全域路由。
  \item[\texttt{features}] 載入選用的功能模組。
\end{description}

\texttt{mainfonts} 可作為 \texttt{globalfonts} 的別名；兩者都表示明確的全域
覆寫；\texttt{fonts} 則提供一般自動介面。

範本介面提供一個可預期的位置，讓完整設定集中表達。包裝類別本身已經套用各自的
版面與全域字體預設，因此文件可以單獨指定額外功能清單。

\subsection{單一用途的捷徑}

若要針對個別子系統進行設定，IMPE 提供對應的公開捷徑：

\begin{verbatim}
\UseLayout{en_doc}
\UseLayouts{...}

\UseFont{libertinus}
\UseFonts{arabic,tibetan}

\UseFont{libertinus}[global]
\UseLocalFonts{arabic,tibetan} % 明確局部覆寫
\UseGlobalFonts{libertinus}    % 明確全域覆寫

\UseFeature{hyperlinks}
\UseFeatures{math,hyperlinks,tables}
\end{verbatim}

對一般使用而言，省略模式的字體形式是建議的捷徑。局部與全域別名則適合在文件刻意
覆寫某個字體族的已登錄行為時使用。統一介面與單一用途介面是同一公開模型的兩種視角，
兩者共用相同的設定語義。

\subsection{可組合的範本}

IMPE 的範本由可重用的元件組成。某個版面可以與不同字體選擇重複使用，
同一組功能可以用於文章或書籍，加入某個書寫系統專用字體族時，其他文件
功能仍維持原有設定。這種組合方式是公開介面的核心設計原則。

例如，下列兩組宣告透過相同範本介面組合出兩份相當不同的文件：

\begin{verbatim}
% 英文文章：類別已提供 en_doc 和 CMU
\documentclass{impeart}
\UseTemplateSet{
  fonts = {libertinus},
  features = {math,hyperlinks}
}

% 中文書籍：類別已提供 zh_book、CMU 與 Shanggu
\documentclass[lang=zh]{impebook}
\UseTemplateSet{
  fonts = {sanskrit,tibetan},
  features = {citations,index,hyperlinks}
}
\end{verbatim}

以上是兩份不同文件的導言區示例。文件類別會提供相應的版面與基礎字體，
作者只需補上額外字體和功能。如需使用其他已登錄版面，再明確指定
\texttt{layout=}。


\section{字體、書寫系統路由與多語排版}

IMPE 提供兩種公開字體選擇方式：簡單的已登錄字體預設，
以及可以自行組合的書寫設定檔與字體資源。兩者都屬 1.1.0
正式支援的介面；既有的字體縮寫命令亦繼續保留。

\subsection{已登錄字體預設}

一般情況下，使用 \verb|\UseFont{id}| 載入既有字體預設；
\verb|\UseTemplateSet| 的 \texttt{fonts} 鍵可以一次指定多個。
省略模式時，IMPE 依完整的登錄規則自動處理：建立局部字體命令、
登錄全域候選字體，或接管指定 Unicode 範圍。

\begin{verbatim}
\UseFont{libertinus}
\UseFonts{sanskrit,tibetan,arabic}
\UseFont{libertinus}[local]
\UseFont{libertinus}[global]
\end{verbatim}

第一個命令既提供 \verb|\LIB|，也把 Libertinus
登錄為使用者明確指定的全域拉丁字體候選。
因此，它可以取代文件類別預設的 CMU。
\texttt{[local]} 只建立局部設定；
\texttt{[global]} 明確要求全域模式。
\verb|\UseLocalFont|、\verb|\UseGlobalFont|
及各自的複數形式，可用來明確指定模式；
\verb|\UseMainFont| 是全域選擇的另一名稱。

像 \verb|\SA{...}| 這類既有命令，
可在文字參數內套用相應字體及專用排版處理。
也可以直接寫 \verb|\Font{sanskrit}{...}|，
不必預先呼叫 \verb|\UseFont|。
全部 59 個原有字體預設及其局部命令列於附錄。

\subsection{書寫設定檔與字體資源}

字體系統區分四類登錄資訊：
\emph{Unicode 區塊資料}、
\emph{書寫設定檔}（碼位覆蓋範圍與文字處理規則）、
\emph{字體資源／字體組}（實際字面與樣式），
以及\emph{已登錄字體預設}。
例如，\texttt{sanskrit} 是包含梵語排版行為的舊式字體預設，
\texttt{devanagari} 則是可以搭配不同字體資源的書寫設定檔。

1.1.0 提供 69 個書寫設定檔、85 個字體資源別名，
並保留原有 59 個字體預設。書寫設定檔決定
Unicode 範圍及文字處理方式，字體資源決定字形來源；
更換字體不會改變該設定檔的碼位範圍。
\texttt{chinese}、\texttt{japanese} 與
\texttt{korean} 的漢字範圍互有重合，
需要時應明確指定相應的地區字形。

\begin{verbatim}
\UseWriting{[shanggu]chinese}
\UseWriting{[tiro]devanagari}
\UseWriting{[naskh]arabic}
\end{verbatim}

一般寫法為 \verb|\UseWriting{[字體資源,選項]設定檔}|，
其中括號內第一項可以省略，也可以寫已登錄的字體資源 ID；
\texttt{font=} 用於指定可解析的字體族名稱，
\texttt{file=} 用於指定字體檔案。這幾種字體來源的選擇
在同一項設定內互斥。多項設定可在括號和群組之外用逗號分隔。
其他可用選項包括 \texttt{stylemap=native|full}、
\texttt{writing=horizontal|vertical}、\texttt{offset=}
及 \texttt{italic=zhuque} 等字面覆寫。
全域字體路由使用橫排；有支援的局部命令則可使用直排。

\subsection{具名與臨時局部字體}

\verb|\DeclareWriting{名稱}{設定}| 可儲存一組設定而不立即啟用；
\verb|\Font{名稱}[覆寫]{文字}| 將其用於指定文字；
\verb|\Writing{設定}{文字}| 則直接使用一組未命名的局部設定。

\begin{verbatim}
\DeclareWriting{han-note}
  {[shanggu,stylemap=full,italic=zhuque]chinese}
\begin{document}
\Font{han-note}{漢字}
\Font{sanskrit}{संस्कृतम्}
\Writing{[tiro]devanagari}{संस्कृतम्}
\end{document}
\end{verbatim}

局部字體選擇會將一套字體及文字處理方式套用到
\emph{整個文字參數}，不會在參數內再次按 Unicode
範圍自動分拆。內層明確指定的局部字體可以覆寫外層選擇。
沒有特殊文字處理需求時，也可只用
\texttt{font=} 或 \texttt{file=} 作局部設定。

\texttt{stylemap=native} 優先使用字體真正提供的樣式字面；
\texttt{stylemap=full} 還可以使用已登錄的跨字體樣式代用。
部分西文字體組預設使用後者。
樣式缺失與某個字形缺失是兩種不同問題。
完整選項可參閱 \texttt{docs/USEWRITING-1.0.5.md}；
附錄另列出全部設定檔與字體資源識別碼。

\subsection{全域優先級與缺字後備}

標準類別預設 CMU；選用 \texttt{lang=zh} 時另有 Shanggu。
\emph{使用者明確指定的全域字體優先於類別預設}。
在同一優先級內，若多套拉丁字體競爭主字體位置，
通常由最先登錄者取得。中文、日文、韓文的
共用範圍則遵循中文先於日文、日文先於韓文的既有規則。

\begin{verbatim}
\UseGlobalFonts{cmu,libertinus}
\end{verbatim}

這個例子裏，CMU 是第一套明確指定的拉丁主字體，
後面的 Libertinus 不會自動取代它；
但 \verb|\UseFont{libertinus}| 可以優先於類別隱含的 CMU。

自動逐字形缺字後備目前用於 CJK 字體。
IMPE 透過 xeCJK 檢查「襯線／無襯線」、
「常規／粗體」、「正體／斜體」所組成的
$2\times2\times2$ 樣式矩陣；
先在目前字體資源中尋找，再查找同一書寫設定檔、
最後查找其他 CJK 設定檔的已登錄資源。
使用者明確登錄的字體排在類別隱含字體之前；
後續載入的資源亦會更新後備鏈。

這與字體名稱無法解析時採用的整套代用字體有所不同。
其他文字系統保留 OpenType 對連續文字的塑形處理，
不執行逐碼位的自動缺字替換。

\subsection{Unicode 範圍與 CJK 路由}

按書寫系統啟用的全域字體可以接管指定 Unicode 區塊，
其餘文字仍沿用原有字體。局部命令在其參數內優先，
結束後恢復原有全域路由。

CJK 使用 xeCJK 的主要、無襯線與等寬字體通道。
共用漢字保留文件已選定的 CJK 字體，
假名與諺文可各自選擇字體範圍。
若要顯示日文或韓文的漢字字形，
可使用 \verb|\JP{...}|、\verb|\KR{...}|
等局部命令。重合的全域設定檔仍遵循既有的路由優先級。

\subsection{字形塑造、斷行與書寫方向}

印度系文字、阿拉伯文等複雜文字的 OpenType 字形塑造，
由 \XeTeX{} 和所用字體負責。
IMPE 提供相應的文字與語言設定，以及梵語天城文的音節斷行、
藏文按音節分隔點斷行、泰文依詞典斷行等處理。
需要由右至左書寫的文字也有對應設定。
蒙古文、滿文和回鶻文支援局部旋轉直排，
與整頁書寫方向的變更不同。

\subsection{XSR 1.3：聖書字、契丹小字與圖片字形}

IMPE 隨附 XSR 1.3，提供埃及聖書字組合排版、
契丹小字堆疊與行內圖片字形功能。
相關文字由各自的字體預設啟用；圖片字形
則由獨立的 \texttt{glyphs} 功能載入。

\begin{description}
  \item[埃及聖書字]
    先載入 \verb|\UseFont{egyptian}|，
    再使用對應的 Unicode 組合控制字元或
    \verb|\IMPEEgyptianText{...}|。
  \item[契丹小字]
    載入 \verb|\UseFont{khitan_small}| 後，
    使用 \verb|\KHS{...}| 進行原生 XSR 堆疊。
    亦保留 \verb|\KHSstack| 與 \verb|\KHSstackblock|。
    原有字簇分隔符與 U+16FE4 填充符會保留；
    堆疊高度不同時會在文字行內對齊。
  \item[圖片字形]
    使用 \verb|\UseFeature{glyphs}| 載入功能，
    以 \verb|\IMPEGlyphRegister| 登錄圖檔，
    再以 \verb|\IMPEGlyph| 插入文字行。
\end{description}

動態渲染需要 Python 執行環境、
XeLaTeX Shell Escape 或 XSR 預處理，
以及可讀取的字體與圖檔。
XSR 可透過 PDF 的 \texttt{/ActualText}
保留聖書字組合和契丹小字字簇原本的線性 Unicode 序列；
實際複製結果仍視 PDF 閱讀器而定。

\subsection{更多字體文件}

字體登錄、樣式對應、文字路由與特殊文字處理詳見
\texttt{docs/FONTS.md} 及 \texttt{docs/USEWRITING-1.0.5.md}。
展示冊收錄各種文字的實際排版效果。

\section{版面與功能}

\subsection{各類出版物的版面預設}

版面預設統一決定紙張、版心、行距、頁眉與章節排版。
IMPE 提供可直接使用的出版版面，作者需要細部修改時，
仍可使用一般 \LaTeX{} 套件。

\begin{description}
  \item[\texttt{en\_doc}、\texttt{zh\_doc}]
    A4 文章版面；左右各 29\,mm、上方 28\,mm、
    下方 32\,mm。
  \item[\texttt{en\_report}、\texttt{zh\_report}]
    獨立的 A4 報告版面；左 31\,mm、右 29\,mm、
    上 29\,mm、下 34\,mm，並有專用頁眉。
  \item[\texttt{en\_book}、\texttt{zh\_book}]
    ISO B5（176$\times$250\,mm）雙面書籍版面；
    內側 19\,mm、外側 31\,mm，裝訂補償 3\,mm。
  \item[\texttt{beamer}、\texttt{beamer\_zh}]
    Madrid 主題的 4:3 簡報版面。
\end{description}

文章與報告都有自己的 A4 版心和文字設定。
另提供 \texttt{beamer\_clean} 與
\texttt{beamer\_clean\_zh}，供需要簡潔型簡報的人選用。
要製作寬螢幕簡報，可設定原生類別選項
\texttt{aspectratio=169}。

\subsection{原生版面調整與局部設定}

標準文章、報告和書籍類別已載入
\texttt{geometry} 與 \texttt{fancyhdr}。
若要精確調整頁邊距或頁眉，建議直接使用這兩套
\LaTeX{} 套件的命令：

\begin{verbatim}
\documentclass{impereport}
\geometry{a4paper,left=30mm,right=30mm}
\fancyhead[R]{頁眉短題}
\end{verbatim}

IMPE 也提供 \verb|\LayoutSetup| 和 \verb|\LayoutSet|
設定已登錄的版面參數；
\texttt{LayoutScope} 可在正文的一段範圍內使用局部設定。

\begin{verbatim}
\documentclass[lang=zh]{impeart}
\LayoutSetup{page.inner=25mm,page.top=24mm}
\begin{document}
\begin{LayoutScope}{text.par-indent=2em}
本段採用局部設定的段首縮排。
\end{LayoutScope}
\end{document}
\end{verbatim}

進階排版可繼續使用原生的頁面、頁眉與 Beamer 命令。
需要重現早期 1.0.5 的版面參數時，可查閱
第~\ref{sec:compatibility} 節的相容設定。

\subsection{功能模組}

各功能模組負責載入配套套件，使用者則可直接使用
這些套件的一般 \LaTeX{} 命令：

\begin{description}
  \item[\texttt{math}] 數學套件、定理環境與數學字體設定。
  \item[\texttt{hyperlinks}] 超連結、PDF 書籤和穩定的跳轉目標。
  \item[\texttt{citations}] 使用 \texttt{biblatex}／\texttt{biber}
    的引文及參考書目。
  \item[\texttt{index}] 使用 \texttt{imakeidx}／\texttt{xindy} 的原生索引。
  \item[\texttt{tables}] 表格、\texttt{booktabs}、
    \texttt{longtable} 及 \texttt{tabularx}。
  \item[\texttt{image}] 圖片、圖說與子圖。
  \item[\texttt{drawing}] TikZ／PGF、\texttt{pgfplots}
    （相容版本 1.18）和 \texttt{forest}。
  \item[\texttt{glyphs}] 以 XSR 登錄並插入行內圖片字形。
  \item[\texttt{headers}] 使用 \texttt{fancyhdr} 設定頁眉。
\end{description}

\begin{verbatim}
\UseFeatures{math,hyperlinks,tables,image,drawing,glyphs}
\end{verbatim}

同一份文件可組合多項功能，重複載入同一功能時不會再次初始化。

\subsection{數學與文獻引用}

\texttt{math} 提供 \texttt{theorem}、\texttt{lemma}、
\texttt{proposition}、\texttt{corollary}、
\texttt{definition}、\texttt{example}、
\texttt{remark} 等定理類環境。
若未另行指定，數學字體仍使用 Computer Modern；
要使用 Libertinus Math，可在載入功能前設定：

\begin{verbatim}
\UseMathFont{libertinus}
\UseFeature{math}
\end{verbatim}

\texttt{citations} 使用原生 \texttt{biblatex}，
因此可照常使用 \verb|\addbibresource|、
\verb|\textcite|、\verb|\parencite|、
\verb|\printbibliography| 等命令。
\verb|\UseCitationStyle| 可以設定 APA、GB/T 7714、
數字式及作者—年份式等引文樣式。

\subsection{以原生命令建立索引}

新文件建議直接使用 \texttt{imakeidx} 的
\verb|\makeindex|、\verb|\index|、
\verb|\printindex|。若由 IMPE 建立預設索引，
每一個一般 \verb|\index| 標記都會有獨立的 PDF
跳轉目標；即使同一頁有多個標記，或文件曾重設頁碼，
也能分別連到對應位置。

\begin{verbatim}
\UseFeature{index}
\begin{document}
佛教\index{佛教}與大乘\index{佛教!大乘}。
後文仍可再次標記佛教\index{佛教}。
\printindex
\end{document}
\end{verbatim}

若使用者已自行宣告預設索引，IMPE 會遵循其設定。
也可以使用原生語法建立具名索引，例如
\verb|\makeindex[name=people,title={人物},program=xindy]|，
搭配 \verb|\index[people]{...}| 和
\verb|\printindex[people]|。
逐項獨立跳轉的保證範圍，是 IMPE 自行建立的
預設索引中的普通索引標記；
\texttt{see}、\texttt{seealso}、範圍標記、
具名索引等進階語法仍遵循原生處理方式。

IMPE 所建立的預設索引可在輸出目錄使用：

\begin{verbatim}
texindy -L english -C utf8 -M <job>-impe <job>.idx
\end{verbatim}

完成索引處理後，再執行兩次 \XeLaTeX{}。

\subsection{表格、圖片、繪圖與行內字形}

一般表格使用 \texttt{table}、\texttt{tabular}、
\texttt{longtable}、\texttt{booktabs} 等原生環境與命令；
圖片使用 \texttt{figure}、\verb|\includegraphics| 和子圖功能；
繪圖則使用 \texttt{tikzpicture}、\texttt{axis}、
\texttt{forest}。舊式的表格與圖片便利命令仍供既有文件使用。

圖片字形的用法如下：

\begin{verbatim}
\UseFeature{glyphs}
\IMPEGlyphRegister[logical={未編碼字形},description={拓片}]
  {bs-042}{images/bs-042.svg}
\begin{document}
前文 \IMPEGlyph{bs-042} 後文。
\IMPEGlyph[scale=1.2,raise=1pt,trim={1pt 0pt 1pt 0pt}]{bs-042}
\end{document}
\end{verbatim}

圖片路徑必須可以由渲染程式讀取。
\texttt{logical} 與 \texttt{description} 用於描述字形；
\texttt{trim} 裁切的是顯示範圍。

\subsection{頁眉與超連結}

\texttt{headers} 可以取 \verb|\title| 的第一行
作為固定頁眉短題，也可以使用
\verb|\HeaderTitle{...}| 另行指定。
更細緻的頁眉變動可直接使用 \texttt{fancyhdr}。

\texttt{hyperlinks} 載入 \texttt{hyperref} 與
\texttt{bookmark}，提供支援 Unicode 的 PDF 設定、
穩定的跳轉目標與可點選的目錄。
文件依然使用原生 \verb|\label|、\verb|\ref|、
\verb|\href|、\verb|\url| 命令。
詳見 \texttt{docs/FEATURES.md}。

\section{相容性與遷移}\label{sec:compatibility}

新文件統一使用 \texttt{impe.sty}、
\texttt{impeart}、\texttt{impereport}、
\texttt{impebook}、\texttt{impebeamer}，
中文以 \texttt{lang=zh} 選擇設定。
既有的 \texttt{impeart\_zh}、\texttt{impereport\_zh}、
\texttt{impebook\_zh}、\texttt{impebeamer\_zh}
分別等同於對應的標準類別加上 \texttt{lang=zh}；
\texttt{next*} 舊類別及 \texttt{nextsystem.sty}
亦繼續作為相容入口。

原有 \verb|\UseFont| 識別碼、局部字體縮寫及
\texttt{[local]}／\texttt{[global]} 模式仍受支援；
新的書寫設定檔與字體資源介面提供另一種組合方式。

索引以 \verb|\index| 等原生命令為推薦方式。
原有 \verb|\Term| 和 \verb|\IndexTitle|
保留其格式與首次出現規則；
一般 \verb|\index| 則能對每次出現的位置建立標記。
需要真正的術語表時，可以直接使用
\texttt{glossaries-extra}。

舊式 \texttt{TableInlineFit}、\texttt{TableLong}、
\texttt{TableBook}、\texttt{TableBookX}、
\texttt{TableBookNotes}、\texttt{NiceBooktable*}、
\texttt{OneImage}、\texttt{OneImageInline}、
\texttt{PanelFigure}、\verb|\Panel| 等介面均保留。
新文件優先使用 \LaTeX{} 的原生表格、圖片與子圖語法；
\texttt{ExampleBlock} 也仍可供舊稿使用。

若要還原 1.0.5 的版面參數，可在標準類別之後使用
\verb|\UseLegacyLayout|：

\begin{verbatim}
\documentclass[lang=zh]{impebook}
\UseLegacyLayout
\end{verbatim}

舊式的 \texttt{report} 及 \texttt{legacy\_*}
版面識別碼仍可使用。相容性保留的是公開命令與設定；
字體或版式更新後，分頁可能有所改變。

直接安裝程式能辨識 \texttt{tex/latex/nextsystem/}
下的 0.1.x 系列安裝，並將目前套件放到
\texttt{tex/latex/impe/}。遷移時會保留既有的使用者檔案，
在適當情況下轉移本機設定，並清理已識別的舊套件檔案。
新的本機設定檔應使用 \texttt{impe.local.tex}。

\section{套件結構與擴充點}

程式碼倉庫將公開入口、穩定機制、登錄資料、特殊實作與本機資源分開：

\begin{verbatim}
package/   公開套件與類別入口
core/      穩定的子系統機制
catalog/   公開識別碼與登錄資料
modules/   可擴充或專門化的實作
assets/    字體等本機執行資源
scripts/   安裝與發行工具
manual/    手冊原始碼
docs/      詳細子系統文件
\end{verbatim}

\subsection{穩定核心}

\texttt{core/} 層包含供新增字體族、版面預設或功能使用的穩定共用機制。它負責
實作字體登錄與路由機制、版面登錄與相容性檢查、功能載入器，以及其他共同行為。

這項區分對維護十分重要。如果加入一個普通 OpenType 字體就必須修改字體路由引擎，
或加入一項功能就必須修改通用功能載入器，代表擴充模型已經失敗。因此，穩定核心刻意
與不斷增長的具體資源目錄分離。

\subsection{登錄表}

\texttt{catalog/} 層負責分配公開識別碼與中繼資料，其中集中保存字體、版面與功能的
主要登錄資料。對很多擴充工作而言，這是最合適的起點。

一個採用標準 \texttt{fontspec} 塑形的普通字體族，通常可以用登錄表中的一筆
註冊來表示。類似地，公開版面預設由現有內部組件組成，而功能識別碼則對應到其實作
模組。將這些宣告集中起來，可以在一處審查公開介面。

\subsection{模組}

\texttt{modules/} 層保留給真正需要獨立實作的功能：功能模組、特殊書寫系統支援，
或可重用的內部版面組件。因此，新模組代表真正的行為擴充；字體名稱則寫入登錄表。

這條原則對多語支援尤其重要。多數書寫系統都應使用一般 OpenType 塑形與登錄表
中繼資料；專門模組服務需要超出通用核心行為的書寫系統或渲染策略。

\subsection{擴充 IMPE}

希望擴充系統的開發者，應先閱讀詳細登錄介面的
\texttt{docs/SYSTEM-zh.md}，再查閱相應的子系統文件：
\texttt{docs/FONTS-zh.md}、\texttt{docs/LAYOUTS-zh.md} 或
\texttt{docs/FEATURES-zh.md}。

一般原則很簡單：登錄資料是一般擴充的首選；特殊行為交由專門模組處理，
讓穩定核心維持共用機制。

\section{進一步文件與回饋}

本手冊是一般使用 IMPE 的起點。程式碼倉庫中還包含更詳細的技術文件，供需要完整參考
或希望擴充系統的使用者查閱。

\begin{description}
  \item[\texttt{docs/SYSTEM-zh.md}] 架構、公開入口層、範本設定、發行模型與擴充邊界。
  \item[\texttt{docs/FONTS-zh.md}] 字體族登錄、局部與全域模式、路由、後備、
    塑形、書寫行為、直排支援與外部渲染。
  \item[\texttt{docs/LAYOUTS-zh.md}] 公開版面預設、相容目標，以及預設所採用的
    組件模型。
  \item[\texttt{docs/FEATURES-zh.md}] 公開功能識別碼，以及各模組提供的命令、
    環境、套件組合與編譯要求。
  \item[\texttt{CHANGELOG-zh.md}] 已發行版本的變更記錄。
  \item[展示冊] 完整視覺展示冊可由
    \href{\IMPEShowcaseURL}{專案程式碼倉庫}取得。
\end{description}

專案程式碼倉庫位於

\begin{center}
\href{\IMPERepo}{\texttt{github.com/KelvinYangBK67/IMPE-LaTeX-System}}
\end{center}

錯誤報告、功能請求、文件修正、可攜性報告以及其他具體回饋，都歡迎透過 GitHub Issues
提交：

\begin{center}
\href{\IMPEIssues}{\texttt{github.com/KelvinYangBK67/IMPE-LaTeX-System/issues}}
\end{center}

若回報排版問題，最好附上一份最小原始碼、所用的 \TeX{} 引擎與 \TeX{} Live 版本、
IMPE 版本，以及相關字體資訊。若問題涉及本機或私下提供的字體，也應盡可能區分
「字體本身缺少字形」與「IMPE 路由錯誤」這兩類情況。

英文手冊是權威手冊文本。其他語言版本可以依各自的技術書寫習慣調整行文與術語，
但其中的命令、範例以及對 IMPE 行為的描述，都應與英文版保持同步。

\appendix

\section{快速參考}
\label{app:quick-reference}

本附錄概述 IMPE 1.x 已記錄的公開介面，可與正文中的說明及 \texttt{docs/} 下更詳細的
子系統文件配合查閱。內部實作命令，尤其是以 \texttt{\string\Next...} 開頭的命令，
由其實作文件說明。

\subsection{建議的字體載入模型}

一般使用時，請優先在 \texttt{\string\UseTemplateSet} 中使用
\texttt{fonts = \{...\}}，或直接使用 \texttt{\string\UseFont\{...\}}。省略
模式時，IMPE 會遵循字體族已登錄的行為：視實際設定而定，它可能建立局部命令、
啟用文件全域字體，或安裝針對特定文字或範圍的自動路由。

明確指定模式可依文件需求刻意覆寫自動行為：

\begin{verbatim}
\UseFont{sanskrit}          % 建議：使用已登錄的自動行為
\UseFont{sanskrit}[local]   % 明確只使用局部模式
\UseFont{sanskrit}[global]  % 明確使用全域／範圍全域模式

\UseTemplateSet{
  fonts = {sanskrit,tibetan,arabic},
  features = {math,hyperlinks}
}
\end{verbatim}

便利命令 \texttt{\string\UseGlobalFont(s)} 與
\texttt{\string\UseLocalFont(s)} 表達的是同一種明確區分。
\texttt{globalfonts} 與 \texttt{mainfonts} 範本鍵同樣屬於明確全域覆寫，
\texttt{fonts} 仍提供一般自動選擇。


\subsection{文件類別與套件入口}

\begin{longtable}{@{}p{0.31\linewidth}p{0.63\linewidth}@{}}
\toprule
\textbf{入口} & \textbf{用途} \\
\midrule
\endfirsthead
\toprule
\textbf{入口} & \textbf{用途} \\
\midrule
\endhead
\texttt{impeart} &
英文文章包裝類別；預設使用 \texttt{en\_doc}，並以 CMU 作為基礎全域字體。 \\
\texttt{impeart\_zh} &
中文文章包裝類別；預設使用 \texttt{zh\_doc}，並載入 CMU 與 Shanggu。 \\
\texttt{impereport} &
英文報告包裝類別；預設使用 \texttt{en\_doc}。 \\
\texttt{impereport\_zh} &
中文報告包裝類別；預設使用 \texttt{zh\_doc}。 \\
\texttt{impebook} &
英文書籍包裝類別；預設使用 \texttt{en\_book}。 \\
\texttt{impebook\_zh} &
中文書籍包裝類別；預設使用 \texttt{zh\_book}。 \\
\texttt{impebeamer} &
英文 Beamer 包裝類別；預設使用 \texttt{beamer} 版面。 \\
\texttt{impebeamer\_zh} &
中文 Beamer 包裝類別；預設使用 \texttt{beamer} 版面，並載入 CMU 與 Shanggu。 \\
\texttt{impe.sty} &
可在文件選用的一般 \LaTeX{} 類別下直接載入 IMPE。 \\
\bottomrule
\end{longtable}

舊的 \texttt{next*} 套件與類別名稱在 1.x 中仍保留為相容別名，但新文件應使用
標準的 \texttt{impe*} 入口。

\subsection{核心設定命令}

\begin{longtable}{@{}p{0.38\linewidth}p{0.56\linewidth}@{}}
\toprule
\textbf{命令} & \textbf{用途} \\
\midrule
\endfirsthead
\toprule
\textbf{命令} & \textbf{用途} \\
\midrule
\endhead
\texttt{\string\UseTemplateSet\{...\}} &
在一個宣告中設定版面、字體與功能。 \\
\texttt{\string\UseFont\{id\}[mode]} &
載入一個已登錄字體族；省略 \texttt{mode} 時採用建議的自動行為，也可明確指定
\texttt{local} 或 \texttt{global}。 \\
\texttt{\string\UseFonts\{ids\}[mode]} &
以相同的選用模式載入多個已登錄字體族。 \\
\texttt{\string\UseGlobalFont\{id\}} &
明確要求一個字體族使用全域模式。 \\
\texttt{\string\UseGlobalFonts\{ids\}} &
明確要求多個字體族使用全域模式。 \\
\texttt{\string\UseMainFont\{id\}} &
\texttt{\string\UseGlobalFont} 的別名；保留作為面向文件設定的名稱。 \\
\texttt{\string\UseMainFonts\{ids\}} &
\texttt{\string\UseGlobalFonts} 的別名。 \\
\texttt{\string\UseLocalFont\{id\}} &
明確要求一個字體族使用局部模式。 \\
\texttt{\string\UseLocalFonts\{ids\}} &
明確要求多個字體族使用局部模式。 \\
\texttt{\string\UseLayout\{id\}} &
載入一個已登錄版面預設。 \\
\texttt{\string\UseLayouts\{ids\}} &
載入多個已登錄版面預設。 \\
\texttt{\string\UseFeature\{id\}} &
載入一個選用功能模組。 \\
\texttt{\string\UseFeatures\{ids\}} &
載入多個選用功能模組。 \\
\texttt{\string\SetCatalogFontRoot\{path\}} &
覆寫字體登錄表解析字體檔時使用的根目錄。 \\
\texttt{\string\subtitle\{text\}} &
設定 IMPE 標題區塊使用的副標題。 \\
\texttt{\string\HeaderTitle\{text\}} &
覆寫 \texttt{headers} 功能使用的固定頁眉標題。 \\
\texttt{\string\HeaderStyle\{style\}} &
選擇頁眉樣式；已記錄的值為 \texttt{running} 與 \texttt{title}。 \\
\bottomrule
\end{longtable}

\subsection{範本組設定鍵}

\begin{longtable}{@{}p{0.23\linewidth}p{0.69\linewidth}@{}}
\toprule
\textbf{鍵} & \textbf{意義} \\
\midrule
\endfirsthead
\toprule
\textbf{鍵} & \textbf{意義} \\
\midrule
\endhead
\texttt{layout} &
選擇一個公開版面預設。 \\
\texttt{fonts} &
依照已登錄的自動行為載入一個或多個字體族；這是一般字體載入的建議設定鍵。 \\
\texttt{globalfonts} &
明確強制列出的字體族使用其全域模式或全域路由。 \\
\texttt{mainfonts} &
\texttt{globalfonts} 的別名。 \\
\texttt{features} &
載入一個或多個選用功能模組。 \\
\bottomrule
\end{longtable}

\subsection{公開版面識別碼}

\begin{longtable}{@{}p{0.23\linewidth}p{0.69\linewidth}@{}}
\toprule
\textbf{ID} & \textbf{用途} \\
\midrule
\endfirsthead
\toprule
\textbf{ID} & \textbf{用途} \\
\midrule
\endhead
\texttt{en\_doc} &
採用 A4 頁面幾何與英文文字間距的英文文章／報告版面。 \\
\texttt{zh\_doc} &
採用 A4 頁面幾何與中文文字間距的中文文章／報告版面。 \\
\texttt{report} &
只適用於 report、並使用與 \texttt{zh\_doc} 相同頁面與正文 component 的相容 preset。 \\
\texttt{en\_book} &
採用書籍頁面幾何、頁眉與右頁開章的英文書籍／報告版面。 \\
\texttt{zh\_book} &
採用書籍頁面幾何、中文文字間距、頁眉與右頁開章的中文書籍／報告版面。 \\
\texttt{beamer} &
使用 IMPE Beamer 字體排印與導覽設定的簡報版面。 \\
\bottomrule
\end{longtable}

\subsection{公開功能識別碼}

\begin{longtable}{@{}p{0.23\linewidth}p{0.69\linewidth}@{}}
\toprule
\textbf{ID} & \textbf{用途} \\
\midrule
\endfirsthead
\toprule
\textbf{ID} & \textbf{用途} \\
\midrule
\endhead
\texttt{math} &
載入 IMPE 數學套件組合與定理環境。 \\
\texttt{hyperlinks} &
載入超連結、PDF 書籤、可點擊標題與雙向腳註連結。 \\
\texttt{citations} &
載入 \texttt{biblatex}/\texttt{csquotes} 引用支援與 IMPE 引用樣式輔助功能。 \\
\texttt{index} &
載入索引支援與 \texttt{\string\Term} 輔助命令。 \\
\texttt{tables} &
載入表格套件組合、欄類型與 IMPE 表格環境。 \\
\texttt{image} &
載入圖片、圖說與多面板圖形輔助功能。 \\
\texttt{lists\_envs} &
載入 \texttt{ExampleBlock} 顯示環境。 \\
\texttt{headers} &
為文章、報告與書籍類文件載入可設定的頁眉。 \\
\texttt{bib}、\texttt{header} &
分別是 \texttt{citations} 與 \texttt{headers} 的相容別名。 \\
\bottomrule
\end{longtable}

\subsection{功能專屬命令與設定值}

\begin{longtable}{@{}p{0.51\linewidth}p{0.43\linewidth}@{}}
\toprule
\textbf{命令或設定} & \textbf{用途／已記錄的值} \\
\midrule
\endfirsthead
\toprule
\textbf{命令或設定} & \textbf{用途／已記錄的值} \\
\midrule
\endhead
\texttt{\string\UseMathFont\{value\}} &
在載入 \texttt{math} 前選擇數學字體路徑：
\texttt{auto}、\texttt{libertinus}、\texttt{newcm}、\texttt{mlmodern}，
或其他可傳給 \texttt{\string\setmathfont} 的字體名稱。 \\
\texttt{\string\UseCitationStyle\{value\}} &
在載入 \texttt{citations} 前選擇引用預設：
\texttt{APA}、\texttt{GB}、\texttt{numeric} 或 \texttt{author-year}。 \\
\texttt{\string\SetCitationBiblatexOptions\{...\}} &
在載入 \texttt{citations} 前取代實際使用的 \texttt{biblatex} 選項清單。 \\
\texttt{\string\HeaderTitle\{text\}} &
設定頁眉使用的簡短固定標題。 \\
\texttt{\string\HeaderStyle\{value\}} &
選擇 \texttt{running} 或 \texttt{title}。 \\
\texttt{\string\IndexTitle} &
可選的索引標題覆寫；需在載入 \texttt{index} 前定義。 \\
\texttt{\string\Term[options]\{display\}[description]} &
以粗體排出術語並索引其第一次出現；選項包括
\texttt{sort}、\texttt{key} 與 \texttt{parentheses}。 \\
\texttt{\string\TablesSetup} &
套用或重新套用 IMPE 表格設定。 \\
\texttt{\string\Panel} &
在 \texttt{PanelFigure} 或 \texttt{PanelFigure*} 中插入一個面板。 \\
\texttt{\string\TemplateFigurePaths} &
設定圖片功能使用的預設搜尋路徑。 \\
\texttt{\string\OneImageDefaultWidth} &
設定 \texttt{OneImage} 的預設寬度。 \\
\texttt{\string\OneImageMaxHeight} &
設定 \texttt{OneImage} 的最大預設圖片高度。 \\
\texttt{\string\OneImageDefaultPlacement} &
設定 \texttt{OneImage} 的預設浮動體位置選項。 \\
\texttt{\string\PanelDefaultCols} &
設定多面板圖形預設使用的欄數。 \\
\texttt{\string\PanelDefaultHeight} &
設定預設面板高度。 \\
\texttt{\string\PanelDefaultMode} &
設定預設面板版面模式。 \\
\texttt{\string\PanelDefaultPlacement} &
設定多面板圖形的預設浮動體位置選項。 \\
\bottomrule
\end{longtable}

由功能所載入套件直接提供的標準命令，例如
\texttt{\string\addbibresource}、\texttt{\string\printbibliography}、
\texttt{\string\printindex}，或 \texttt{booktabs} 的橫線命令，仍使用原套件的
正常 \LaTeX{} 語法。

\subsection{公開環境與表格欄類型}

\begin{longtable}{@{}p{0.37\linewidth}p{0.57\linewidth}@{}}
\toprule
\textbf{名稱} & \textbf{用途} \\
\midrule
\endfirsthead
\toprule
\textbf{名稱} & \textbf{用途} \\
\midrule
\endhead
\texttt{theorem}, \texttt{lemma}, \texttt{proposition}, \texttt{corollary} &
由 \texttt{math} 功能提供的帶編號定理類環境。 \\
\texttt{definition}, \texttt{example}, \texttt{remark} &
由 \texttt{math} 功能提供的其他定理樣式環境。 \\
\texttt{TableInlineFit} &
將緊湊表格配合可用寬度排版的環境。 \\
\texttt{TableLong} &
供內容跨頁使用的長表格環境。 \\
\texttt{TableBook}, \texttt{TableBookX} &
書籍樣式的表格環境。 \\
\texttt{TableBookNotes} &
帶附註的書籍樣式表格環境。 \\
\texttt{NiceBooktable}, \texttt{NiceBooktableX}, \texttt{NiceBooktableNotes} &
由 \texttt{tables} 功能提供的另一組書籍樣式表格輔助環境。 \\
\texttt{L}, \texttt{C}, \texttt{R} &
分別為右側不齊、置中與左側不齊的 \texttt{tabularx} 欄。 \\
\texttt{P\{w\}}, \texttt{M\{w\}}, \texttt{B\{w\}} &
分別採右側不齊、置中與左側不齊的固定寬度段落欄。 \\
\texttt{OneImage} &
單圖圖形環境；在 Beamer 下改為行內內容而非浮動體。 \\
\texttt{OneImageInline} &
置中的非浮動圖片。 \\
\texttt{PanelFigure} &
帶總圖說與子圖說的多面板圖形；在一般文件類別下使用。 \\
\texttt{PanelFigure*} &
無圖說的多面板圖形。 \\
\texttt{ExampleBlock} &
帶縮排的斜體顯示區塊，適合範例、引文、語言材料或教學內容。 \\
\bottomrule
\end{longtable}

\subsection{已登錄字體族識別碼}

下表列出 IMPE 1.x 中可傳給 \texttt{\string\UseFont} 的字體族 ID。
\textbf{自動}欄概括省略模式的 \texttt{\string\UseFont\{id\}} 一般會啟用的行為：

\begin{description}
  \item[\texttt{G}] 文件全域字體族；
  \item[\texttt{L}] 局部字體族／命令；
  \item[\texttt{L+G}] 同時提供局部存取與完整文件全域啟用；
  \item[\texttt{L+R}] 同時提供局部存取與依文字或範圍限制的自動全域路由。
\end{description}

自動路徑是建議使用的介面。明確的 \texttt{[local]} 或 \texttt{[global]} 會刻意繞過
部分自動判斷，適合刻意覆寫時使用。

\begin{longtable}{@{}p{0.30\linewidth}p{0.11\linewidth}p{0.18\linewidth}p{0.31\linewidth}@{}}
\toprule
\textbf{字體族 ID} & \textbf{自動} & \textbf{局部命令} & \textbf{典型用途} \\
\midrule
\endfirsthead
\toprule
\textbf{字體族 ID} & \textbf{自動} & \textbf{局部命令} & \textbf{典型用途} \\
\midrule
\endhead
\texttt{cmu} & G & --- & Computer Modern Unicode 基礎字體族。 \\
\texttt{noto} & L+G & \texttt{\string\NOT} & 一般 Noto 拉丁字體族。 \\
\texttt{times} & L+G & \texttt{\string\TIM} & Times／Arial／Consolas 系統字體族路徑。 \\
\texttt{gentium} & L+G & \texttt{\string\GEN} & Gentium Plus。 \\
\texttt{charis} & L+G & \texttt{\string\CHA} & Charis SIL。 \\
\texttt{libertinus} & L+G & \texttt{\string\LIB} & Libertinus 襯線／無襯線／等寬字體族。 \\
\texttt{mlmodern} & L & \texttt{\string\MLM} & Latin Modern／MLModern 局部路徑。 \\

\texttt{anatolian} & L & \texttt{\string\CA} & 卡里亞文／安納托利亞文字材料。 \\
\texttt{coptic} & L & \texttt{\string\CO} & 科普特文。 \\
\texttt{bopomofo} & L & \texttt{\string\ZY} & 注音符號。 \\
\texttt{cuneiform} & L & \texttt{\string\CU} & 楔形文字。 \\
\texttt{glagolitic} & L & \texttt{\string\GL} & 格拉哥里字母。 \\
\texttt{italic} & L & \texttt{\string\OI} & 古意大利字母。 \\
\texttt{hungarian} & L & \texttt{\string\OH} & 古匈牙利字母。 \\
\texttt{runic} & L & \texttt{\string\RU} & 盧恩文字。 \\
\texttt{armenian} & L & \texttt{\string\HY} & 亞美尼亞文。 \\
\texttt{georgian} & L & \texttt{\string\KA} & 格魯吉亞文。 \\

\texttt{hindi} & L+R & \texttt{\string\HI} & 帶已登錄範圍路由的印地語天城文。 \\
\texttt{sanskrit} & L+R & \texttt{\string\SA} & 帶梵語專用路由的天城文。 \\
\texttt{devanagari} & L & \texttt{\string\DEV} & 通用天城文塑形。 \\
\texttt{tamil} & L & \texttt{\string\TA} & 泰米爾文。 \\
\texttt{brahmi} & L & \texttt{\string\BR} & 婆羅米文。 \\
\texttt{tibetan} & L+R & \texttt{\string\TI} & 帶有按音節點感知斷行行為的藏文。 \\
\texttt{segoe} & L & \texttt{\string\SEG} & Segoe Historic 局部路徑。 \\
\texttt{thai} & L & \texttt{\string\TH} & 泰文。 \\

\texttt{arabic} & L & \texttt{\string\AR} & 阿拉伯文；提供襯線、無襯線與替代樣式字面。 \\
\texttt{urdu} & L & \texttt{\string\UR} & 烏爾都語納斯塔利克體。 \\
\texttt{aramaic} & L & \texttt{\string\IA} & 帝國阿拉姆文。 \\
\texttt{nabataean} & L & \texttt{\string\NB} & 納巴泰文。 \\
\texttt{hebrew} & L & \texttt{\string\HE} & 希伯來文。 \\
\texttt{syriac} & L & \texttt{\string\SY} & 敘利亞文。 \\
\texttt{syriac\_eastern} & L+R & \texttt{\string\SYE} & 帶已登錄範圍路由的東部敘利亞文。 \\
\texttt{kharosthi} & L & \texttt{\string\KH} & 佉盧文。 \\
\texttt{khitan\_small} & L & \texttt{\string\KHS} & 使用專用組合器的契丹小字。 \\
\texttt{pahlavi\_parthian} & L+R & \texttt{\string\PAR} & 銘文帕提亞文。 \\
\texttt{pahlavi\_inscriptional} & L+R & \texttt{\string\PAH} & 銘文巴列維文。 \\
\texttt{pahlavi\_psalter} & L & \texttt{\string\PSP} & 帶專門支援的詩篇巴列維文。 \\
\texttt{avestan} & L & \texttt{\string\AV} & 阿維斯陀文。 \\
\texttt{manichaean} & L & \texttt{\string\MA} & 摩尼文。 \\
\texttt{phoenician} & L & \texttt{\string\PH} & 腓尼基文。 \\
\texttt{samaritan} & L & \texttt{\string\SM} & 撒馬利亞文。 \\
\texttt{sogdian} & L & \texttt{\string\SG} & 粟特文。 \\
\texttt{sogdian\_old} & L+R & \texttt{\string\SGO} & 帶已登錄範圍路由的古粟特文。 \\

\texttt{chinese\_simplified} & L+R & \texttt{\string\SC} & 簡體中文 CJK 字體族。 \\
\texttt{chinese\_traditional} & L+R & \texttt{\string\TC} & 繁體中文 CJK 字體族。 \\
\texttt{japanese} & L+R & \texttt{\string\JP} & 日文字體族；自動路由保留共用漢字的優先次序。 \\
\texttt{wenjin} & L+R & \texttt{\string\WJ} & 帶 CJK 後備鏈的 WenJin Mincho。 \\
\texttt{shanggu} & G & --- & 面向繁體中文的全域 CJK 字體族。 \\
\texttt{sim} & G & --- & 面向簡體中文的系統 CJK 字體族。 \\
\texttt{korean} & L+R & \texttt{\string\KR} & 韓文／諺文字體族。 \\
\texttt{tangut} & L+R & \texttt{\string\TG} & 帶 CJK 感知自動路由的西夏文。 \\

\texttt{mongolian} & L & \texttt{\string\MO} & 帶內建直排版面路徑的蒙古文。 \\
\texttt{mongolian\_baiti} & L & \texttt{\string\MOb} & 帶直排版面路徑的 Mongolian Baiti。 \\
\texttt{manchu} & L & \texttt{\string\MC} & 帶直排版面路徑的滿文。 \\
\texttt{turkic} & L & \texttt{\string\OT} & 古突厥文。 \\
\texttt{uyghur} & L & \texttt{\string\UY} & 帶直排／旋轉版面路徑的古維吾爾文。 \\

\texttt{vietnamese\_quocngu} & L & \texttt{\string\VI} & 越南國語字。 \\
\texttt{vietnamese\_hannom} & L & \texttt{\string\HN} & 越南漢喃文／CJK 文字。 \\
\bottomrule
\end{longtable}

\subsection{進階擴充介面}

下列命令供擴充 IMPE 使用。

\begin{longtable}{@{}p{0.39\linewidth}p{0.55\linewidth}@{}}
\toprule
\textbf{命令} & \textbf{用途} \\
\midrule
\endfirsthead
\toprule
\textbf{命令} & \textbf{用途} \\
\midrule
\endhead
\texttt{\string\FontRegisterFamily\{...\}} &
註冊新的字體族及其局部／全域中繼資料。 \\
\texttt{\string\LayoutPresetRegister\{...\}} &
利用內部版面組件註冊一個結構化版面預設。 \\
\texttt{\string\LayoutPresetDeclare\{...\}} &
建置系統時使用的較低階版面預設宣告介面。 \\
\bottomrule
\end{longtable}

完整的登錄語法與擴充規則，請參閱 \texttt{docs/FONTS-zh.md}、
\texttt{docs/LAYOUTS-zh.md} 與 \texttt{docs/SYSTEM-zh.md}。

\section{展示冊}\label{app:showcase}

以下各頁完整收入標準 IMPE 展示冊。展示冊是參考手冊的視覺配套文件，
用來展示多語字體路由、複雜塑形、由右至左文字、直排、CJK 地區字形，以及各種文件
功能的組合效果。

最新展示冊可由
\href{\IMPEShowcaseURL}{目前主分支}取得；與本手冊對應的版本則由
\href{\IMPETaggedShowcaseURL}{v\IMPEVersion{} 標籤}保存。

\IfFileExists{\IMPEShowcaseFile}
  {\includepdf[pages=1-3,pagecommand={}]{\IMPEShowcaseFile}%
   \includepdf[pages=4-6,landscape=true,pagecommand={}]{\IMPEShowcaseFile}%
   \includepdf[pages=7-,pagecommand={}]{\IMPEShowcaseFile}}
  {\PackageError{impe-manual}
     {缺少必要的程式碼倉庫資源 ../showcase/impe-showcase.pdf}
     {請從 manual/zh-tw 編譯，或在程式碼倉庫根目錄執行 scripts/build_manual.ps1。}}

\end{document}
